\chapter{Geometri Koordinat}\label{ch:g10:coordgeom}

Gagasan besar Descartes adalah memerikan titik dengan pasangan bilangan,
sehingga soal geometri berubah menjadi perhitungan. Hanya dengan dua rumus
--- \omterm{prop:g10:coordgeom:midpoint}{titik tengah} dan jarak --- orang dapat membuktikan bahwa sebuah segitiga
sama kaki, bahwa sebuah segi empat adalah jajargenjang, atau bahwa tiga titik
segaris, tanpa menarik satu pun garis bantu.

\section{Koordinat pada bidang}

\begin{definition}[Sistem koordinat]\label{def:g10:coordgeom:system}
Sebuah \emph{sistem koordinat}\index{sistem koordinat} pada bidang terdiri
atas sebuah titik asal $O$ dan dua sumbu berskala yang melalui $O$: sumbu
$x$ yang mendatar dan sumbu $y$ yang tegak. Setiap titik $M$ lalu mempunyai
sepasang \emph{koordinat}\index{koordinat} $(x, y)$ yang tunggal:
\emph{absis}\index{absis} $x$ dan \emph{ordinat}\index{ordinat} $y$. Sistem
itu disebut \emph{ortonormal}\index{sistem ortonormal} apabila kedua sumbunya
tegak lurus dan membawa satuan panjang yang sama. Semua sistem dalam bab
ini bersifat ortonormal.
\end{definition}

\begin{omfigure}
\begin{tikzpicture}[scale=0.85]
  \draw[thin, gray!40] (-2.4,-1.4) grid (4.4,3.4);
  \draw[-Stealth] (-2.6,0) -- (4.6,0) node[below, font=\small] {$x$};
  \draw[-Stealth] (0,-1.6) -- (0,3.6) node[left, font=\small] {$y$};
  \node[below left, font=\small] at (0,0) {$O$};
  \fill[omDef] (3,2) circle (2pt);
  \node[above right, font=\small, omDef] at (3,2) {$M(3, 2)$};
  \draw[dashed, omDef] (3,0) -- (3,2) -- (0,2);
  \fill[omThm] (-2,-1) circle (2pt);
  \node[below left, font=\small, omThm] at (-2,-1) {$N(-2, -1)$};
  \foreach \x in {1,2,3,4} \node[below, font=\footnotesize] at (\x,0) {\x};
  \foreach \y in {1,2,3} \node[left, font=\footnotesize] at (0,\y) {\y};
\end{tikzpicture}

{\small Setiap titik pada bidang ditentukan oleh dua bilangan: absisnya
dahulu (mendatar), lalu \omterm{def:g10:coordgeom:system}{ordinatnya} (tegak).}
\end{omfigure}

\section{Titik tengah sebuah ruas garis}

\begin{proposition}[Rumus titik tengah]\label{prop:g10:coordgeom:midpoint}
Misalkan $A(x_A, y_A)$ dan $B(x_B, y_B)$. \emph{Titik tengah}\index{titik tengah}
$I$ ruas garis $[AB]$ berkoordinat
\[
I\left(\frac{x_A + x_B}{2},\ \frac{y_A + y_B}{2}\right).
\]
\end{proposition}

\begin{proof}
Tinjaulah \omterm{def:g10:coordgeom:system}{koordinat} mendatarnya. Dari $A$ ke $B$, absisnya
berubah sebesar $x_B - x_A$; titik di tengah jalan berabsis
\[
x_A + \frac{x_B - x_A}{2} = \frac{2x_A + x_B - x_A}{2}
= \frac{x_A + x_B}{2},
\]
yaitu rata-rata kedua absisnya. Perhitungan yang sama berlaku untuk
\omterm{def:g10:coordgeom:system}{ordinatnya}.
\end{proof}

\begin{example}\label{ex:g10:coordgeom:midpoint}
Dengan $A(-1, 4)$ dan $B(5, -2)$, \omterm{prop:g10:coordgeom:midpoint}{titik tengah} $[AB]$ adalah
\[
I\left(\frac{-1 + 5}{2},\ \frac{4 + (-2)}{2}\right) = I(2, 1).
\]
Rumus itu berjalan mundur juga: jika $A(-1, 4)$ dan \omterm{prop:g10:coordgeom:midpoint}{titik tengahnya}
$I(2, 1)$, maka $B$ memenuhi $\frac{-1 + x_B}{2} = 2$ dan
$\frac{4 + y_B}{2} = 1$, jadi $x_B = 5$ dan $y_B = -2$: titik $B$ adalah
\emph{pencerminan} $A$ terhadap $I$.
\end{example}

\section{Jarak antara dua titik}

\begin{theorem}[Rumus jarak]\label{thm:g10:coordgeom:distance}
Pada \omterm{def:g10:coordgeom:system}{sistem ortonormal}, jarak antara $A(x_A, y_A)$ dan
$B(x_B, y_B)$ adalah
\[
AB = \sqrt{(x_B - x_A)^2 + (y_B - y_A)^2}.
\]
\end{theorem}

\begin{proof}
Misalkan $C$ titik $(x_B, y_A)$: tingginya sama dengan $A$ dan
absisnya sama dengan $B$, sehingga segitiga $ACB$ bersudut siku-siku di $C$.
Kedua sisi sikunya adalah ruas garis mendatar dan tegak, dengan panjang
$AC = \abs{x_B - x_A}$ dan $CB = \abs{y_B - y_A}$. Menurut teorema
Pythagoras,
\[
AB^2 = AC^2 + CB^2 = (x_B - x_A)^2 + (y_B - y_A)^2,
\]
lalu menarik akar kuadratnya (kedua ruas taknegatif) memberi rumus tadi.
\end{proof}

\begin{omfigure}
\begin{tikzpicture}[scale=0.85]
  \draw[thin, gray!40] (-0.4,-0.4) grid (5.4,3.4);
  \draw[-Stealth] (-0.6,0) -- (5.6,0) node[below, font=\small] {$x$};
  \draw[-Stealth] (0,-0.6) -- (0,3.6) node[left, font=\small] {$y$};
  \coordinate (A) at (1,1);
  \coordinate (B) at (5,3);
  \coordinate (C) at (5,1);
  \draw[omDef, thick] (A) -- (B);
  \draw[omThm, thick] (A) -- (C) node[midway, below, font=\small]
    {$\abs{x_B - x_A}$};
  \draw[omThm, thick] (C) -- (B) node[midway, right, font=\small]
    {$\abs{y_B - y_A}$};
  \draw[thin] (C) ++(-0.25,0) -- ++(0,0.25) -- ++(0.25,0);
  \fill (A) circle (2pt) node[above left, font=\small] {$A$};
  \fill (B) circle (2pt) node[above, font=\small] {$B$};
  \fill (C) circle (2pt) node[below right, font=\small] {$C$};
\end{tikzpicture}

{\small Rumus jarak adalah teorema Pythagoras yang diterapkan pada
segitiga siku-siku $ACB$ yang dibangun pada sisi mendatar dan sisi tegak.}
\end{omfigure}

\begin{example}\label{ex:g10:coordgeom:distance}
Dengan $A(1, 1)$ dan $B(5, 3)$:
\[
AB = \sqrt{(5-1)^2 + (3-1)^2} = \sqrt{16 + 4} = \sqrt{20} = 2\sqrt 5 .
\]
Jarak biasanya dibiarkan dalam bentuk eksak (dengan akar); bulatkan hanya
di akhir sekali bila jawaban desimal memang diperlukan.
\end{example}

\begin{remark}
Urutan titiknya tidak berpengaruh: $(x_A - x_B)^2 =
(x_B - x_A)^2$. Namun jangan lupakan kuadratnya! Jaraknya
\emph{bukan} $\abs{x_B - x_A} + \abs{y_B - y_A}$.
\end{remark}

\section{Memakai kedua rumus dalam geometri}

\begin{method}[Jenis sebuah segitiga]\label{met:g10:coordgeom:triangle}
Diberikan tiga titik $A$, $B$, $C$ lewat \omterm{def:g10:coordgeom:system}{koordinatnya}:
\begin{enumerate}
  \item hitunglah ketiga kuadrat panjangnya $AB^2$, $AC^2$, $BC^2$ dengan
        rumus jarak (mempertahankan kuadratnya menghindarkan akar);
  \item dua kuadrat panjang yang sama $\Rightarrow$ segitiganya
        \emph{sama kaki};
  \item jika kuadrat panjang terbesarnya sama dengan jumlah dua yang lain,
        segitiga itu \emph{siku-siku} di titik sudut yang berhadapan dengan
        sisi terpanjangnya, menurut kebalikan teorema Pythagoras.
\end{enumerate}
\end{method}

\begin{example}\label{ex:g10:coordgeom:triangle}
Misalkan $A(0, 1)$, $B(4, 3)$ dan $C(2, -3)$. Maka
\begin{align*}
AB^2 &= (4-0)^2 + (3-1)^2 = 16 + 4 = 20, \\
AC^2 &= (2-0)^2 + (-3-1)^2 = 4 + 16 = 20, \\
BC^2 &= (2-4)^2 + (-3-3)^2 = 4 + 36 = 40 .
\end{align*}
Pertama, $AB^2 = AC^2$, jadi $AB = AC$: segitiga itu sama kaki di $A$.
Lebih lanjut $AB^2 + AC^2 = 40 = BC^2$, jadi menurut kebalikan teorema
Pythagoras segitiga itu juga siku-siku di $A$.
\end{example}

\begin{method}[Mengenali jajargenjang]\label{met:g10:coordgeom:parallelogram}
Sebuah segi empat $ABCD$ adalah jajargenjang tepat ketika kedua diagonalnya
$[AC]$ dan $[BD]$ mempunyai \omterm{prop:g10:coordgeom:midpoint}{titik tengah} yang sama. Jadi: hitunglah kedua
\omterm{prop:g10:coordgeom:midpoint}{titik tengah} itu dengan rumus \omterm{prop:g10:coordgeom:midpoint}{titik tengah} lalu bandingkan.
\end{method}

\begin{example}\label{ex:g10:coordgeom:parallelogram}
Misalkan $A(-1, 0)$, $B(2, 2)$, $C(5, 1)$ dan $D(2, -1)$. \omterm{prop:g10:coordgeom:midpoint}{Titik tengah}
$[AC]$ adalah $\left(\frac{-1+5}{2}, \frac{0+1}{2}\right) = (2, \frac12)$;
\omterm{prop:g10:coordgeom:midpoint}{titik tengah} $[BD]$ adalah $\left(\frac{2+2}{2}, \frac{2-1}{2}\right) =
(2, \frac12)$. Keduanya berimpit, jadi $ABCD$ adalah jajargenjang.
\end{example}

\begin{omfigure}
\begin{tikzpicture}[scale=0.75]
  \draw[thin, gray!40] (-1.4,-1.4) grid (5.4,2.4);
  \draw[-Stealth] (-1.6,0) -- (5.6,0);
  \draw[-Stealth] (0,-1.6) -- (0,2.6);
  \coordinate (A) at (-1,0);
  \coordinate (B) at (2,2);
  \coordinate (C) at (5,1);
  \coordinate (D) at (2,-1);
  \draw[omDef, thick] (A) -- (B) -- (C) -- (D) -- cycle;
  \draw[omThm, dashed] (A) -- (C);
  \draw[omThm, dashed] (B) -- (D);
  \fill (A) circle (2pt) node[left, font=\small] {$A$};
  \fill (B) circle (2pt) node[above, font=\small] {$B$};
  \fill (C) circle (2pt) node[right, font=\small] {$C$};
  \fill (D) circle (2pt) node[below, font=\small] {$D$};
  \fill[omThm] (2,0.5) circle (2pt) node[above right, font=\small] {$I$};
\end{tikzpicture}

{\small $ABCD$ adalah jajargenjang karena kedua diagonalnya berbagi \omterm{prop:g10:coordgeom:midpoint}{titik
tengah} yang sama, yaitu $I(2, \frac12)$.}
\end{omfigure}

\section{Latihan}

\begin{exercise}[$\star$]\label{exo:g10:coordgeom:1}
Misalkan $A(2, 5)$ dan $B(-4, 1)$. Hitunglah \omterm{def:g10:coordgeom:system}{koordinat} \omterm{prop:g10:coordgeom:midpoint}{titik tengah}
$[AB]$ dan jarak $AB$.
\end{exercise}

\begin{exercise}[$\star$]\label{exo:g10:coordgeom:2}
Misalkan $A(3, -2)$ dan $I(1, 2)$. Carilah \omterm{def:g10:coordgeom:system}{koordinat} titik $B$ sehingga
$I$ menjadi \omterm{prop:g10:coordgeom:midpoint}{titik tengah} $[AB]$.
\end{exercise}

\begin{exercise}[$\star$]\label{exo:g10:coordgeom:3}
Di antara titik $P(4, 1)$, $Q(-3, 2)$ dan $R(0, -5)$, manakah yang terdekat
ke titik asal $O(0,0)$? Jawablah dengan jarak eksak.
\end{exercise}

\begin{exercise}[$\star$]\label{exo:g10:coordgeom:4}
Gambarlah titik $A(1, 2)$, $B(5, 4)$, $C(7, 0)$ pada kertas berpetak, lalu
tunjukkan lewat perhitungan bahwa segitiga $ABC$ sama kaki. Apakah siku-siku?
\end{exercise}

\begin{exercise}[$\star\star$]\label{exo:g10:coordgeom:5}
Misalkan $A(-2, 1)$, $B(1, 3)$, $C(4, 1)$ dan $D(1, -1)$.
\begin{enumerate}
  \item Tunjukkan bahwa $ABCD$ adalah jajargenjang.
  \item Hitunglah $AB$ dan $BC$. Apakah $ABCD$ belah ketupat (semua sisinya sama)?
  \item Hitunglah $AC$ dan $BD$. Apakah $ABCD$ persegi panjang (diagonalnya sama)?
\end{enumerate}
\end{exercise}

\begin{exercise}[$\star\star$]\label{exo:g10:coordgeom:6}
Lingkaran berpusat $\Omega(2, 1)$ dan berjari-jari $5$ terdiri atas semua
titik yang berjarak $5$ dari $\Omega$. Di antara titik $A(5, 5)$,
$B(-1, -3)$, $C(6, 2)$, manakah yang terletak pada lingkaran itu? Di dalamnya? Di luarnya?
\end{exercise}

\begin{exercise}[$\star\star$]\label{exo:g10:coordgeom:7}
Misalkan $A(1, 1)$ dan $B(7, 5)$. Carilah semua titik $M(x, 0)$ pada sumbu
$x$ yang berjarak sama dari $A$ dan $B$. (Tulislah $MA^2 = MB^2$ lalu
selesaikan untuk $x$.)
\end{exercise}

\begin{exercise}[$\star\star$]\label{exo:g10:coordgeom:8}
Misalkan $A(0, 3)$, $B(4, 1)$.
\begin{enumerate}
  \item Hitunglah \omterm{def:g10:coordgeom:system}{koordinat} \omterm{prop:g10:coordgeom:midpoint}{titik tengah} $I$ dari $[AB]$.
  \item Tunjukkan bahwa $M(2, 2) = I$, lalu periksa dengan menghitung $MA$
        dan $MB$ bahwa $M$ berjarak sama dari $A$ dan $B$.
  \item Carilah titik kedua, pada sumbu $y$, yang berjarak sama dari $A$ dan
        $B$.
\end{enumerate}
\end{exercise}

\begin{exercise}[$\star\star$]\label{exo:g10:coordgeom:9}
Diberikan titik $A(-1, -1)$, $B(3, 1)$ dan $C(11, 5)$. Hitunglah $AB$,
$BC$ dan $AC$, lalu simpulkan bahwa $A$, $B$, $C$ segaris. (Tiga titik
segaris tepat ketika jarak terbesar di antara ketiganya sama dengan jumlah
dua jarak yang lain.)
\end{exercise}

\begin{exercise}[$\star\star\star$]\label{exo:g10:coordgeom:10}
Misalkan $A(a, 0)$ dan $B(0, b)$ dengan $a, b > 0$, dan misalkan $I$ \omterm{prop:g10:coordgeom:midpoint}{titik
tengah} $[AB]$. Tunjukkan lewat perhitungan bahwa $OI = IA = IB$, dengan $O$
titik asalnya. (Ini membuktikan sebuah teorema klasik: pada segitiga
siku-siku, \omterm{prop:g10:coordgeom:midpoint}{titik tengah} sisi miringnya berjarak sama dari ketiga titik sudut.)
\end{exercise}

\section{Soal: Jembatan Descartes --- persamaan lingkaran dan cara
mengetahui posisimu}

\begin{problem}[{Soal akhir pekan --- lingkaran menjadi sebuah
persamaan, teorema lama menjadi perhitungan, dan tiga suar
menentukan letak sebuah titik: trilaterasi secara aljabar}]
\label{pb:g10:coordgeom:1}
Pada tahun 1637 René Descartes membangun jembatan antara dua benua:
setiap kurva geometri menjadi sebuah \emph{\omterm{def:g10:algebra:equation}{persamaan}}, setiap
pertanyaan geometri menjadi perhitungan. Soal ini menyusuri jembatan itu
ke dua arah --- menurunkan \omterm{def:g10:algebra:equation}{persamaan} lingkaran dari
rumus jarak (\cref{thm:g10:coordgeom:distance}), membuktikan ulang
teorema klasik dalam tiga baris aljabar, lalu berakhir di tempat
jembatan itu paling ramai dilalui hari ini: menghitung posisi
dari isyarat jarak, yaitu jantung datar sistem penentu posisi satelit.

\medskip
\noindent\textbf{Bagian I --- Lingkaran memperoleh \omterm{def:g10:algebra:equation}{persamaan}.}
\begin{enumerate}
  \item Sebuah titik $M(x, y)$ terletak pada lingkaran berpusat
        $\Omega(2, 1)$ dan berjari-jari $5$ tepat ketika
        $\Omega M = 5$. Kuadratkan syarat ini untuk memperoleh
        \omterm{def:g10:algebra:equation}{persamaan} lingkarannya:
        \[
        (x - 2)^2 + (y - 1)^2 = 25 .
        \]
  \item Ujilah keanggotaannya: di antara $A(5, 5)$, $B(6, 4)$,
        $D(4, 4)$, manakah yang berada pada lingkaran, di dalamnya, di luarnya?
  \item Sebuah lingkaran dalam samaran: lengkapkan kuadratnya
        (\cref{pb:g10:algebra:1}) pada
        \[
        x^2 + y^2 - 6x + 4y - 12 = 0
        \]
        lalu berikan pusat dan jari-jarinya.
  \item Lakukan hal yang sama pada $x^2 + y^2 - 6x + 4y + 14 = 0$. Apa
        yang disingkapkan bentuk lengkapnya? Nyatakan kriterianya: kapan
        $(x - h)^2 + (y - k)^2 = c$ memerikan sebuah lingkaran, sebuah
        titik tunggal, atau bukan apa-apa?
  \item Potonglah lingkaran pada pertanyaan 3 dengan garis mendatar
        $y = 3$: substitusikan lalu selesaikan. Berapa titik potongnya, dan
        apa nama geometris garis semacam itu?
\end{enumerate}

\medskip
\noindent\textbf{Bagian II --- Teorema lama dalam tiga baris.}
\begin{enumerate}[resume]
  \item Teorema lingkaran di jilid sebelumnya (sudut siku-siku pada setiap
        setengah lingkaran), dibuktikan ulang secara aljabar: misalkan
        $A(-r, 0)$ dan $B(r, 0)$ kedua ujung sebuah diameter dan $M(x, y)$
        sebarang titik pada lingkaran $x^2 + y^2 = r^2$. Hitunglah
        $MA^2 + MB^2$ lalu simpulkan dengan kebalikan teorema Pythagoras
        bahwa $\widehat{AMB}$ siku-siku.
  \item Teorema garis berat: dengan $A(-a, 0)$, $B(a, 0)$ (sehingga
        \omterm{prop:g10:coordgeom:midpoint}{titik tengah} $[AB]$ adalah titik asal $I$), tunjukkan bahwa untuk
        setiap titik $M(x, y)$:
        \[
        MA^2 + MB^2 = 2\,MI^2 + \frac{AB^2}{2} .
        \]
  \item Simpulkan tempat kedudukan titik $M$ dengan
        $MA^2 + MB^2 = 26$ ketika $AB = 6$: kurva apa, pusatnya di mana,
        jari-jarinya berapa?
  \item Tempat kedudukan yang lain: $A(0,0)$, $B(3,0)$; carilah semua titik
        dengan $MA = 2\,MB$. (Kuadratkan, jabarkan, lengkapkan
        kuadratnya: sebuah lingkaran termasyhur muncul --- Apollonius sudah
        mengenalnya tanpa \omterm{def:g10:coordgeom:system}{koordinat}.)
  \item Dalam satu atau dua kalimat: apa yang diubah jembatan Descartes
        tentang \emph{cara} teorema dibuktikan? Bandingkan
        pertanyaan 6 dengan bukti persegi panjang pada teorema lingkaran
        di jilid sebelumnya (sudut siku-siku pada setiap setengah lingkaran).
\end{enumerate}

\medskip
\noindent\textbf{Bagian III --- Tiga suar menemukanmu.}
Posisimu $M(x, y)$ belum diketahui. Suar $P(0, 0)$ mengukur
jarakmu sebagai $5$; suar $Q(6, 0)$ juga mengukur $5$.
\begin{enumerate}[resume]
  \item Tulislah kedua \omterm{def:g10:algebra:equation}{persamaan} lingkarannya, kurangkan, lalu amati
        kuadratnya saling menghapus: \omterm{def:g10:algebra:equation}{persamaan} sederhana apa yang tersisa?
        Padukan dengan salah satu lingkaran untuk menemukan \emph{dua}
        calon posisimu.
  \item Suar ketiga $R(0, 8)$ mengukur jarakmu sebagai $5$.
        Hitunglah jaraknya ke masing-masing calon lalu putuskan di mana
        kamu berada.
  \item Jelaskan mengapa pengurangan dua \omterm{def:g10:algebra:equation}{persamaan} lingkaran
        \emph{selalu} menghasilkan \omterm{def:g10:algebra:equation}{persamaan} sebuah garis (suku mana
        yang saling menghapus?), dan apa arti geometris garis itu ketika
        kedua lingkarannya berpotongan di dua titik.
  \item Penentuan posisi satelit yang sesungguhnya bekerja di ruang dan
        dengan jam yang tidak sempurna: setiap satelit memberi (lewat waktu
        tempuh isyarat) satu jarak, jadi satu bola. Berapa banyak besaran
        tak diketahui yang dipunyai sebuah penerima (posisinya \emph{dan}
        galat jamnya sendiri), dan mengapa karena itu sistem penentu posisi
        mendengarkan paling sedikit \emph{empat} satelit?
  \item Ketelitian: andaikan jarak suar $P$ sebenarnya
        $5.1$ dan bukan $5$ (yang lain tetap). Ulangi
        pengurangan pada pertanyaan 11 untuk mencari $x$ yang baru, lalu
        ukurlah seberapa jauh taksirannya bergeser. (Bandingkan dengan
        gagasan anggaran galat pada \cref{pb:g10:numbers:1}.)
\end{enumerate}

\medskip
\noindent\textbf{Bagian IV --- Kotak perkakas juru ukur.}
\begin{enumerate}[resume]
  \item Jalur menyiram terpendek: sebuah perkemahan di $A(1, 5)$, sebuah
        lumbung di $B(7, 3)$, sebuah sungai lurus sepanjang sumbu $y = 0$.
        Untuk pergi dari $A$ ke sungai lalu ke $B$ dengan
        langkah paling sedikit: cerminkan $B$ terhadap sungai menjadi $B'$,
        lalu carilah tempat ruas garis $[AB']$ memotong sungai.
        Berikan titik potongnya dan panjang total minimalnya.
        (Gagasannya adalah latihan pantulan bola bilyar di jilid sebelumnya,
        kini dapat dihitung sepenuhnya.)
  \item Buktikan keminimalannya terhadap seorang penantang: hitunglah
        jalur total lewat $Q(3, 0)$ lalu periksa bahwa jalur itu kalah
        oleh jawabanmu.
  \item Golongkan segi empat $A(1,1)$, $B(4,2)$,
        $C(5,5)$, $D(2,4)$ selengkapnya: jajargenjang? belah ketupat?
        persegi panjang? persegi?
        (\cref{met:g10:coordgeom:parallelogram},
        \cref{met:g10:coordgeom:triangle} --- bandingkan diagonalnya
        juga.)
  \item Diberikan $A(-1, 2)$, $B(3, 4)$, $C(6, 0)$, carilah $D$ sehingga
        $ABCD$ menjadi jajargenjang (kiat \omterm{prop:g10:coordgeom:midpoint}{titik tengah} diagonal pada
        soal akhir pekan tentang setengah putaran di jilid sebelumnya, kini
        dalam rumus).
  \item Penutup: tiga perkakas --- rumus \omterm{prop:g10:coordgeom:midpoint}{titik tengah}, rumus
        jarak, dan pengurangan \omterm{def:g10:algebra:equation}{persamaan} lingkaran. Untuk masing-masing,
        sebutkan jenis pertanyaan yang dituntaskannya dalam soal ini, lalu
        nyatakan apa yang ditambahkan dua bab berikutnya ke dalam kotak itu
        (arah dan \omterm{def:g10:reffunc:affine}{gradien}: vektor dan \omterm{def:g10:algebra:equation}{persamaan} garis).

\end{enumerate}
\end{problem}
